\subsection{Постановление Правительства Российской Федерации № 1119}

Постановление Правительства Российской Федерации от 01.11.2012 № 1119 «Об утверждении требований к защите персональных данных при их обработке в информационных системах персональных данных» \cite{PP1119} устанавливает требования к обеспечению безопасности персональных данных при их обработке в информационных системах персональных данных (ИСПДн). Данный нормативный акт конкретизирует положения Федерального закона № 152-ФЗ и определяет подходы к выбору и реализации мер защиты информации.

В соответствии с постановлением вводится понятие уровня защищённости персональных данных, который определяется в зависимости от категории обрабатываемых персональных данных, количества субъектов персональных данных, а также актуальных угроз безопасности. Выделяются четыре уровня защищённости, для каждого из которых устанавливается необходимый состав организационных и технических мер по защите информации.

Для разрабатываемого программного средства данный нормативный акт имеет практическое значение в тех случаях, когда система анализа audit-логов Kubernetes-кластера используется в составе ИСПДн или обрабатывает данные, которые могут быть отнесены к персональным. В такой ситуации необходимо определить актуальный уровень защищённости системы и реализовать соответствующий ему набор мер защиты.

К числу таких мер относятся:
\begin{itemize}
    \item обеспечение контроля доступа в помещения, где размещены информационные системы с ПД, чтобы исключить неконтролируемое проникновение посторонних лиц;
    \item утверждение руководителем оператора документа, в котором чётко прописано, кто и на каком основании имеет доступ к персональным данным;
    \item регистрация и учёт событий безопасности, включая действия пользователей и администраторов;
    \item использование сертифицированных средств защиты информации (СЗИ);
\end{itemize}

Таким образом, Постановление Правительства РФ № 1119 определяет необходимость формального подхода к оценке защищённости системы и выбору мер защиты персональных данных. При внедрении системы обнаружения аномалий необходимо учитывать уровень защищённости ИСПДн и обеспечивать соответствие реализуемых механизмов требованиям данного нормативного акта.